
\documentclass[a4paper,11pt]{article}
\usepackage{enumitem}
\usepackage[T1]{fontenc}


\title{Ein kurzer Artikel über Redundanz}
\author{Stefan Wohlfeil}
\date{19.03.2012}
\begin{document}
\maketitle 
\section{Zusammenfassung:}
In diesem Artikel geht es um:
Lorem ipsum dolor sit amet, consectetuer adipiscing elit.
Aenean commodo ligula eget dolor.
 Aenean massa.
Cum sociis natoque penatibus et magnis dis parturient montes,
nascetur ridiculus mus.
Donec quam felis, ultricies nec, pellentesque eu, pretium quis, sem.
Nulla consequat massa quis enim.
Donec pede justo, fringilla vel, aliquet nec, vulputate eget, arcu.
In enim justo, rhoncus ut, imperdiet a, venenatis vitae, justo.
Nullam dictum felis eu pede mollis pretium.
Integer tincidunt.
Cras dapibus.
Vivamus elementum semper nisi.
Aenean vulputate eleifend \\ \\
\\ \\
\\ \\
\\ \\

\tableofcontents 

\section{Der Redundanzbegriff}

Welche verschiednenen Interpretationen des Begriffs Redundanz gibt es
in welchen Kontexten?
Diese Frage wird in diesem Abschnitt ausführlich diskutiert. \\
\\ \\
\\ \\
\subsection{ Redundanz in der Informatik}

Weit hinten, hinter den Wortbergen, fern der Länder Vokalien und
Konsonantien leben die Blindtexte.
Abgeschieden wohnen Sie in Buchstabhausen an der Küste des
Semantik, eines großen Sprachozeans.
Ein kleines Bächlein namens Duden fließt durch ihren Ort und
versorgt sie mit den nötigen Regelialien.
Es ist ein paradiesmatisches Land, in dem einem gebratene Satzteile
in den Mund fliegen.

Nicht einmal von der allmächtigen Interpunktion werden die
Blindtexte beherrscht – ein geradezu unorthographisches Leben.
Eines Tages aber beschloß eine kleine Zeile Blindtext,
ihr Name war Lorem Ipsum, hinaus zu gehen in die weite Grammatik.
Der große Oxmox riet ihr davon ab, da es dort wimmele von bösen
Kommata, wilden Fragezeichen und hinterhältigen Semikoli, doch
das Blindtextchen ließ sich nicht beirren.

Es packte seine sieben Versalien, schob sich sein Initial in den
Gürtel und machte sich auf den Weg.
Als es die ersten Hügel des Kursivgebirges erklommen hatte,
warf es einen letzten Blick zurück auf die Skyline seiner
Heimatstadt Buchstabhausen, die Headline von Alphabetdorf und
die Subline seiner eigenen Straße, der Zeilengasse.
Wehmütig lief ihm eine rethorische Frage über die Wange, dann
setzte es seinen Weg fort.

Unterwegs traf es eine Copy.
Die Copy warnte das Blindtextchen, da, wo sie herkäme wäre sie
zigmal umgeschrieben worden und alles, was von ihrem Ursprung
noch übrig wäre, sei das Wort "und" und das Blindtextchen solle
umkehren und wieder in sein eigenes, sicheres Land zurückkehren.
Doch alles Gutzureden konnte es nicht überzeugen und so dauerte
es nicht lange, bis ihm ein paar heimtückische Werbetexter
auflauerten, es mit Longe und Parole betrunken machten und
es dann in ihre Agentur schleppten, wo sie es für ihre Projekte
wieder und wieder mißbrauchten.
Und wenn es nicht umgeschrieben wurde, dann benutzen Sie es immernoch.
Weit hinten, hinter den Wortbergen, fern der Länder Vokalien und
Konsonantien leben die Blindtexte.
Abgeschieden\\ 
\\ \\
\\ \\



\subsection{Redundanz in der Philosophie}

Es gibt im Moment in diese Mannschaft, oh, einige Spieler vergessen
ihnen Profi was sie sind.
Ich lese nicht sehr viele Zeitungen, aber ich habe gehört viele Situationen.
Erstens: wir haben nicht offensiv gespielt.
Es gibt keine deutsche Mannschaft spielt offensiv und die Name
offensiv wie Bayern.

Letzte Spiel hatten wir in Platz drei Spitzen: Elber, Jancka und dann Zickler.
Wir müssen nicht vergessen Zickler.
Zickler ist eine Spitzen mehr, Mehmet eh mehr Basler.
Ist klar diese Wörter, ist möglich verstehen, was ich hab gesagt? Danke.
Offensiv, offensiv ist wie machen wir in Platz.

Zweitens: ich habe erklärt mit diese zwei Spieler: nach
Dortmund brauchen vielleicht Halbzeit Pause.
Ich habe auch andere Mannschaften gesehen in Europa nach diese Mittwoch.
Ich habe gesehen auch zwei Tage die Training.
Ein Trainer ist nicht ein Idiot! Ein Trainer sei sehen was passieren in Platz.
In diese Spiel es waren zwei, drei diese Spieler waren schwach wie
eine Flasche leer!
Haben Sie gesehen Mittwoch, welche Mannschaft hat gespielt Mittwoch?
Hat gespielt Mehmet oder gespielt Basler oder hat gespielt Trapattoni?
Diese Spieler beklagen mehr als sie spielen!
Wissen Sie, warum die Italienmannschaften kaufen nicht diese Spieler?
Weil wir haben gesehen viele Male solche Spiel! Haben \\
\\
\\


 \subsection{ Redundanz in der Politik}

Eine wunderbare Heiterkeit hat meine ganze Seele eingenommen, gleich
den süßen Frühlingsmorgen, die ich mit ganzem Herzen genieße.
Ich bin allein und freue mich meines Lebens in dieser Gegend,
die für solche Seelen geschaffen ist wie die meine.
Ich bin so glücklich, mein Bester, so ganz in dem Gefühle von
ruhigem Dasein versunken, daß meine Kunst darunter leidet.

Ich könnte jetzt nicht zeichnen, nicht einen Strich, und bin nie
ein größerer Maler gewesen als in diesen Augenblicken.

Wenn das liebe Tal um mich dampft, und die hohe Sonne an der
Oberfläche der undurchdringlichen Finsternis meines Waldes ruht,
und nur einzelne Strahlen sich in das innere Heiligtum stehlen,
ich dann im hohen Grase am fallenden Bache liege, und näher an der
Erde tausend mannigfaltige Gräschen mir merkwürdig werden; wenn
ich das Wimmeln der kleinen Welt zwischen Halmen, die unzähligen,
unergründlichen Gestalten der Würmchen, der Mückchen näher an meinem
Herzen fühle, und fühle die Gegenwart des Allmächtigen, der uns
nach seinem Bilde schuf, das Wehen des Alliebenden, der uns in
ewiger Wonne schwebend trägt und erhält; mein Freund!
Wenn's dann um meine Augen dämmert, und die Welt um mich her und der
Himmel ganz in meiner Seele ruhn wie die Gestalt einer \\
\\
\\
\\
\\

\section{Praktische Einsatzgebiete der Redundanz} 

\subsection{ Redundanz zur Fehlererkennung in der IT}

Er hörte leise Schritte hinter sich.  Das bedeutete nichts Gutes.
Wer würde ihm schon folgen, spät in der Nacht und dazu noch in
dieser engen Gasse mitten im übel beleumundeten Hafenviertel?

Gerade jetzt, wo er das Ding seines Lebens gedreht hatte und mit der
Beute verschwinden wollte! Hatte einer seiner zahllosen Kollegen
dieselbe Idee gehabt, ihn beobachtet und abgewartet, um ihn nun
um die Früchte seiner Arbeit zu erleichtern?

Oder gehörten die Schritte hinter ihm zu einem der unzähligen
Gesetzeshüter dieser Stadt, und die stählerne Acht um seine
Handgelenke würde gleich zuschnappen?
Er konnte die Aufforderung stehen zu bleiben schon hören.

Gehetzt sah er sich um.
Plötzlich erblickte er den schmalen Durchgang.
Blitzartig drehte er sich nach rechts und verschwand zwischen
den beiden Gebäuden.
Beinahe wäre er dabei über den umgestürzten Mülleimer
gefallen, der mitten im Weg lag.
Er versuchte, sich in der Dunkelheit seinen Weg zu ertasten und \\
\\
\\


\subsection{ Redundanz zur Vertuschung fehlender Inhalte}

Gerade in der Politik fehlen oft Inhalte oder sie sollen dem Wahlvolk
absichtlich vorenthalten werden.
Das verbessert die Wahlchancen.
Jemand musste Josef K.
 verleumdet haben, denn ohne dass er etwas
Böses getan hätte, wurde er eines Morgens verhaftet.
Wie ein Hund! sagte er, es war, als sollte die Scham ihn überleben.
Als Gregor Samsa eines Morgens aus unruhigen Träumen erwachte,
fand er sich in seinem Bett zu einem ungeheueren Ungeziefer verwandelt.

Und es war ihnen wie eine Bestätigung ihrer neuen Träume und guten
Absichten, als am Ziele ihrer Fahrt die Tochter als erste sich erhob
und ihren jungen Körper dehnte.
»Es ist ein eigentümlicher Apparat«, sagte der Offizier zu dem
Forschungsreisenden und überblickte mit einem gewissermaßen
bewundernden Blick den \\
\\
\\
\\
\\

\section{Zusammenfassung und Fazit}

Hier eine kurze Liste der in diesem Dokument vorgestellten
"Arten" von Redundanz:
\begin{enumerate}[label={}]
\item 1. Redundanz in der Informatik
\item 2. Redundanz in der Philosophie
\item 3. Redundanz in der Politik
\end{enumerate}

Zwei flinke Boxer jagen die quirlige Eva und ihren Mops durch Sylt.
Franz jagt im komplett verwahrlosten Taxi quer durch Bayern.
Zwölf Boxkämpfer jagen Viktor quer über den großen Sylter Deich.

Vogel Quax zwickt Johnys Pferd Bim. Sylvia wagt quick den Jux bei Pforzheim.
Polyfon zwitschernd aßen Mäxchens Vögel Rüben, Joghurt und
Quark. "Fix, Schwyz! " quäkt Jürgen blöd vom Paß.

Victor jagt zwölf Boxkämpfer quer über den großen Sylter Deich.
Falsches Üben von Xylophonmusik quält jeden größeren Zwerg.
Heizölrückstoßabdämpfung. Zwei flinke Boxer jagen die quirlige
Eva und ihren Mops durch Sylt.

Franz jagt im komplett verwahrlosten Taxi quer durch Bayern.
Zwölf Boxkämpfer jagen Viktor quer über den großen Sylter Deich.
Vogel Quax zwickt Johnys Pferd Bim. Sylvia wagt quick den Jux bei
Pforzheim. Polyfon zwitschernd aßen Mäxchens Vögel Rüben, Joghurt und
Quark. "Fix, Schwyz! " quäkt Jürgen blöd vom Paß.
Victor jagt zwölf Boxkämpfer quer über den großen Sylter Deich. Falsches 
\end{document}

